\section{Ablauf Unterrichtseinheit (für Lehrpersonen)}
\subsection{Lektion 1: Einführung in Netzwerke}
\begin{todolist}
    \item Einführung in Netzwerke: Was ist ein Netzwerk? Welche Arten von Netzwerken gibt es?
    \item Geschichte des Internets
    \item Einführung in das OSI-Modell
    \item Praktische Übung auf Papier: Nachrichten verschicken
    \item Sicherung: Quiz zu den Themen der Lektion, Moodle
\end{todolist}

\subsection{Lektion 2: Netzwerkprotokolle und -dienste}
\begin{todolist}
    \item Einführung in Netzwerkprotokolle: TCP/IP, HTTP, FTP, DNS
    \item Praktische Übung: Nutzung von Netzwerkdiensten (z.B. Webbrowser, FTP-Client)
    \item Diskussion: Wie funktionieren diese Protokolle im Alltag?
    \item Sicherung: Quiz zu den Protokollen und Diensten, Moodle
\end{todolist}

\subsection{Lektion 3: Netzwerksicherheit}
\begin{todolist}
    \item Einführung in Netzwerksicherheit: Firewalls, Verschlüsselung, VPNs
    \item Praktische Übung: Einrichten einer einfachen Firewall-Regel
    \item Diskussion: Warum ist Netzwerksicherheit wichtig?
    \item Sicherung: Quiz zu Netzwerksicherheit, Moodle
\end{todolist}

\subsection{Lektion 4: Netzwerktechnologien und -architekturen}
\begin{todolist}
    \item Einführung in Netzwerktechnologien: LAN, WAN, WLAN, Ethernet
    \item Praktische Übung: Aufbau eines einfachen Netzwerks mit Routern und Switches
    \item Diskussion: Vor- und Nachteile verschiedener Netzwerktechnologien
    \item Sicherung: Quiz zu Netzwerktechnologien, Moodle
    \item Abschlussdiskussion: Wie beeinflussen Netzwerke unseren Alltag?
\end{todolist}
